\section{Evaluation Methodology}
\label{sec:method}

We measure a real system under a controlled workload~\cite{jain1991art}; factorial designs, randomization, and blocking follow standard experimental design~\cite{montgomery2017design}. The measured host is a VMware virtual machine (8 vCPUs of a Xeon Platinum 8276M, one NUMA node, 16\,GB, vmxnet3) running Linux 7.0 and bpftrace 0.25; some experiments add an identical second one on the same physical host and virtual switch, with an NVIDIA RTX~3060 (12\,GB) serving Ollama 0.40.6.

\subsection{Workloads and Reference}

Synthetic executions run under four conditions: \textbf{N0} distinct process trees and distinct cgroups (control); \textbf{N1} distinct trees sharing one cgroup v2, running concurrently; \textbf{N2} one worker process serving operations of all executions sequentially in random order; and \textbf{N3} one worker serving them concurrently on a pool of 8 threads. Each operation is marked with the runtime context marker and opens a TCP connection, either in the execution's own thread (\emph{inproc}) or through a forked and exec'd subprocess (\emph{subproc}), so that context inheritance across process boundaries is exercised.

Before each run, every operation is assigned an exclusive local destination port, drawn at random and stored in the run plan. The reference origin of a captured connection is obtained from its destination. Resolvers receive only $(t, \mathit{pid}, \mathit{tid}, \mathit{cgroup}, \mathit{ctx})$, and a static check over the resolver code rejects any reference to destination fields. The reference therefore shares no input with the signals under evaluation, \ie it is an oracle independent of the system under test~\cite{barr2015oracle}. Connections to ports outside the plan are counted and excluded; planned operations without a captured connection are reported as capture loss.

The four resolvers (cgroup-only, lineage-only, context-only, and fused) are applied offline to the same event log of each run, so differences between resolvers cannot come from different executions. The same logs feed emulations of the mechanisms of related systems (\cf Section~\ref{sec:results}).

To inject faults, a fraction $p \in \{0, 0.05, 0.2\}$ of each execution's operations is labeled with another execution's identifier, modeling runtime bugs in the context binding, following controlled fault-injection practice~\cite{arlat1990fault,hsueh1997fault}. The number of mislabeled operations is exact ($p \times$ operations), with positions drawn at random, so that fault injection adds no sampling variance.

\subsection{Experiments}
\textbf{E1} (correctness) is a full factorial over structure (N0--N3), $p$, and operation mode, with two executions of 100 operations each. \textbf{E2} (scale) varies the number of concurrent executions ($K \in \{2,4,\dots,64\}$) and the connection rate, with a random pause of up to 10\,ms between connections (\emph{paced}) or none (\emph{max}), in N1 and N3; \textbf{E2b} (capacity) runs $P$ concurrent instances of the closed-loop \texttt{connect} microbenchmark of E3 to find where capture stops being complete. \textbf{E3} (overhead) measures \texttt{connect} latency and throughput with that C microbenchmark, with and without the collector and the context markers, and \textbf{E3b} splits this cost across four configurations: no instrumentation, probes that only count events, the full collector, and an optimized collector. \textbf{E1fast} reruns E1 with the optimized collector to check that its outcomes do not change.

The next experiments use a real Model Context Protocol server, built with FastMCP, the official Python software development kit (SDK), over Server-Sent Events. In \textbf{E4}, $K \in \{2,4\}$ client sessions share it and issue 25 tool calls each; the tool runs synchronously, in a worker thread, as a subprocess, or as an \texttt{async} coroutine, and emits markers at its start and end. \textbf{E5} compares this manual binding with the task-scoped binding on the same server and tools, with $K \in \{2,4,8\}$; in the automatic mode the tools emit no markers and the identity travels in \texttt{\_meta}. \textbf{E6} adds the egress policy. Executions alternate between a permissive class (all planned destinations) and a restricted one (a random half of its own), so a mislabel can cause both false allows and false denies. In \textbf{E7}, four LLM agents (Qwen2.5-3B served locally by Ollama, temperature 0) perform health-check tasks through a shared, instrumented MCP server with HTTP, subprocess, and TCP-check tools, with the policy active; the model, not a script, decides which tools to call. \textbf{E7m} repeats E7 with Qwen2.5-7B and Llama-3.1-8B on the GPU of the second machine and with four and eight agents. \textbf{E7c} and \textbf{E7o} use commercial models from two leading providers, a high-end and an entry-level tool-calling model of each: Anthropic's Claude Opus~5.5 and Haiku~5.5 (Messages API) and OpenAI's GPT-6.1 Sol and GPT-6 Luna (Responses API, reasoning effort \emph{low}). \textbf{E7r} replays every recorded E7c and E7o run without models (same calls, arguments, labels, and timing).

Finally, \textbf{E5f}, \textbf{E6f}, and \textbf{E7f} repeat E5--E7 with the task-scoped binding while a fraction $p \in \{0.04, 0.2\}$ of each execution's tool calls carries another execution's identity (in E7f, one of each agent's first three calls). \textbf{E6L} runs the policy in N0 and N1 with $p \in \{0.05, 0.2\}$ and three lineage modes (label only, deny on conflict, tree owner), and \textbf{E6Ls} checks, in the shared MCP server without faults, that the lineage check never vetoes a correct label.

\textbf{E5net}, \textbf{E6net}, and \textbf{E3net} repeat E5, E6, and E3 with the services on the second machine, so connections leave through the virtual network interface instead of loopback. In \textbf{E5ap}, AgentProvenance~v0.9.0~\cite{byteyellow2026agentprovenance} runs next to \solution{} in the E5 rounds: its eBPF sensor captures the same connections, we register one binding per tool call as its harness integrations do (call, session, the MCP server's process and cgroup, and the call window seen by the client), and its own engine assigns each connection, scored against the same reference.

\subsection{Metrics and Analysis}
For connections with a valid reference: C/N, A/N, M/N, I/N, and C/(C+I) per resolver (\cf Section~\ref{sec:background}). For accounting, each execution's connections are split into correctly credited, credited to a wrong execution, and uncredited (ambiguous or missing). We do not report the net accounting error alone, because opposite misattributions cancel in it. For E3: median and 99th-percentile latency, throughput, and eBPF nanoseconds per program run.

The unit of repetition is a complete run, with a fresh collector process, new cgroups, and new services; events within a run are not independent replicates~\cite{georges2007statistically}. Run order is randomized with a recorded seed against ordering bias~\cite{mytkowicz2009producing}, and a repetition uses the same seed in every cell (common random numbers). Boots are treated as a higher level of variation~\cite{kalibera2013rigorous}: E1, E2b, E3, E3b, E5--E7, their network, model, and AgentProvenance variants, and the fault-injection experiments are spread over two boots. The sample size depends on the type of claim. For properties expected to hold exactly (I=0 in E5--E7), the primary result is the number of runs with any error; as a secondary summary we give the one-sided 95\% Clopper--Pearson upper bound on the per-connection error rate~\cite{clopper1934use}, about $3/n$ with zero errors in $n$ cases~\cite{hanley1983nothing}, so we planned at least 299 cases per cell (a 1\% bound): 6 runs per cell in E5, E6, and their variants, and 26 per E7, E7m, E7c, and E7o cell, counted on decisions; for AgentProvenance, 300 or more connections per cell give a 95\% interval of $\pm$5.7 points. It assumes independent connections, which does not strictly hold within a run. For performance (E3, E3b, E3net), runs per boot are sized from the pilot variance for a 95\% confidence-interval half-width of 5\%~\cite{hoefler2015scientific}. E1 is deterministic given the labels and is reported as counts. Every campaign used here is \emph{confirmatory}: it ran from a clean revision after a passing preflight, with no failed run; pilots are never results.

\subsection{Reproducibility}
\label{sec:repro}

User tools (Python, R, Snakemake) are pinned by a Nix flake; the kernel, bpftrace, and bpftool are recorded per campaign. A Snakemake workflow and one script per boot run preflight, campaigns, validation, and report. Each campaign is an immutable directory with run order, raw events, plans, scores, environment snapshots, git revision, manifest, and SHA-256 checksums, and every table and figure is generated by script from these directories~\cite{bajpai2019dagstuhl,collberg2016repeatability}.
The code and workflow and summarized results are available at~\cite{sourceCode}.
